\documentclass[11pt]{article}
\usepackage{fontspec,amsmath,amssymb,amsthm,hyperref}
\usepackage[margin=1in]{geometry}
\setlength{\emergencystretch}{3em}
\title{Exact future Brownian moments and derived orthogonality}
\author{LangevinLean mathematical exposition}
\date{}
\begin{document}\maketitle
\noindent Source module: \texttt{AverageLMC/GaussianBrownian.lean}. Compile this proof document with XeLaTeX or LuaLaTeX. The accompanying HTML pairs the mathematical argument with the exact Lean statements and records the scope of the formalization.


The improved energy estimate uses the actual Brownian process underlying the stationary Langevin solution. Its future increments have exact moments, and their orthogonality to the coupled past is derived from independence and adaptedness.



For the scaled noise \(B_t=\sqrt2W_t\), genuine independence and the constructed adapted integral solution imply, for \(0\le t\le h\),


\[\mathbb E(B_h-B_t)=0,\qquad\|B_h-B_t\|_{L^2}^{2}=2(h-t)p.\]
\[\langle X+B_t-Y_t,\,B_h-B_t\rangle_{L^2}=0.\]
\section{Center and measure future Brownian noise}


The Langevin additive noise is √2 times a Brownian increment. Its expectation vanishes at every time, so the expectation of a future difference also vanishes. Algebraically, the difference between the noises at h and t is the additive noise over the segment starting at t with length h−t. The existing exact Brownian L\ensuremath{{}^{2}} norm formula therefore gives its squared norm as 2(h−t)d. This uses the true elapsed-time Gaussian variance and covers the endpoint t=h.



\noindent Lean declarations: \texttt{brownian\_future\_mean\_zero}, \texttt{brownian\_future\_norm\_sq}.

\section{Derive the orthogonality of the actual error}


The complete initial state and stopped Brownian past are independent of the future Brownian increment. Projection onto the initial-state and current-noise coordinates gives zero L\ensuremath{{}^{2}} inner product with each of them because the future increment is centered. For Y(t), the natural-past measurability of the actual integral solution has already been proved by the invariant-subspace construction; applying it to the identity vector field gives the same zero inner product. Subtract and add the three identities to prove that X+B(t)−Y(t) is orthogonal to B(h)−B(t). No stochastic-error orthogonality condition is assumed in the final LMC theorem.



\noindent Lean declarations: \texttt{gaussian\_error\_future\_inner\_zero}.

\end{document}
